% Part: computability
% Chapter: recursive-functions
% Section: non-pr-functions

\documentclass[../../../include/open-logic-section]{subfiles}

\begin{document}

\olfileid{cmp}{rec}{npr}
\olsection{بنیادی بازگشتی نہ ہون والے فنکشن}

بنیادی بازگشتی فنکشن، بدیہی سمجھ مطابق قابلِ حساب
سارے فنکشناں نوں نہیں سمیٹدے۔ بدیہی طور اُتے ایہہ
صاف ہونا چاہیدا اے کہ اسی سارے اک-جگہ والے بنیادی
بازگشتی فنکشناں دی فہرست
$f_0$، $f_1$، $f_2$،~\dots ایویں بنا سکدے آں
کہ $f_x$ دی قیمت انپٹ $y$ اُتے موثر طریقے نال
حساب کر سکّیے؛ دوجے لفظاں وچ، فنکشن $g(x,y)$،
جیہدی تعریف ایہہ اے:
\[
g(x,y) = f_x(y)
\]
قابلِ حساب اے۔ پر فیر ایہہ فنکشن وی قابلِ حساب اے:
\begin{eqnarray*}
h(x) & = & g(x,x) + 1 \\
& = & f_x(x) +1.
\end{eqnarray*}
ہر بنیادی بازگشتی فنکشن $f_i$ لئی، $h$ تے
$f_i$ دی قیمت $i$ اُتے وکھری اے۔ ایس لئی
$h$ قابلِ حساب اے، پر بنیادی بازگشتی نہیں؛ تے
$g$ بارے وی ایہو گل کہی جا سکدی اے۔ ایہہ کانتور
دی قطری دلیل دا اک «موثر» روپ اے۔

قابلِ حساب پر بنیادی بازگشتی نہ ہون والے فنکشناں
دیاں ہور صاف مثالاں وی دتیاں جا سکدیاں نیں۔ مثال
لئی، علامت $g^n(x)$ نوں $g(g(\dots g(x)))$
لئی ورتو، جیہدے وچ کل $n$ وار $g$ آؤندا اے؛
تے فنکشناں دا سلسلہ $g_0,g_1,\dots$ ایویں
متعین کرو:
\begin{eqnarray*}
g_0(x) & = & x+1 \\
g_{n + 1}(x) & = & g_n^x(x)
\end{eqnarray*}
تسی جانچ سکدے او کہ ہر فنکشن $g_n$ بنیادی
بازگشتی اے۔ ہر اگلا فنکشن پچھلے توں بہت تیز
ودھدا اے؛ $g_1(x)$، $2x$ دے برابر اے،
$g_2(x)$، $2^x \cdot x$ دے برابر اے، تے
$g_3(x)$ لگ بھگ $x$ واری $2$ نوں اک دوجے
دی قوت بنا کے اُتے اُتے رکھن والے مینار وانگ
ودھدا اے۔ Ackermann--P\'eter فنکشن بنیادی طور
اُتے $G(x) = g_x(x)$ اے، تے دکھایا جا سکدا اے
کہ ایہہ ہر بنیادی بازگشتی فنکشن توں تیز ودھدا اے۔

ہُن بنیادی بازگشتی فنکشناں دی گنتی وار فہرست
بناون دے مسئلے ول واپس آؤ۔ یاد رکھو کہ ہر
بنیادی بازگشتی فنکشن نوں اسی اک علامتی لکھت
دِتّی اے؛ ایس لئی علامتی لکھتاں دی گنتی وار فہرست
بناونا کافی اے۔ اسی قدرتی عدد $\#(F)$ نوں
ہر علامتی لکھت $F$ نال بازگشتی طور اُتے ایویں
جوڑ سکدے آں:
\begin{eqnarray*}
\#(0) & = & \langle 0 \rangle \\
\#(S) & = & \langle 1 \rangle \\
\#(\Proj{n}{i}) & = & \langle 2, n, i \rangle \\
\#(\fn{Comp}_{k,l}[H,G_0,\dots,G_{k-1}]) & = & \langle
3,k,l,\#(H),\#(G_0),\dots,\#(G_{k-1}) \rangle \\
\#(\fn{Rec}_l[G,H]) & = & \langle 4, l, \#(G), \#(H) \rangle
\end{eqnarray*}
ایتھے اسی ایہہ حقیقت ورت رہے آں کہ پچھلے حصے دے
کوڈاں نال عدداں دے ہر سلسلے نوں قدرتی عدد دے
روپ وچ ویکھیا جا سکدا اے۔ نتیجہ ایہہ اے کہ ہر
کوڈ نوں اک قدرتی عدد ملدا اے۔ بے شک، کجھ سلسلے
(تے ایس لئی کجھ عدد) کسے علامتی لکھت دے مطابق
نہیں ہندے؛ پر اسی $f_i$ نوں اوہ اک-جگہ والا
بنیادی بازگشتی فنکشن رکھ سکدے آں جیہدی علامتی
لکھت دا کوڈ $i$ اے، جے $i$ واقعی ایہو جیہی
لکھت دا کوڈ ہووے؛ نہیں تاں مستقل $0$ فنکشن
رکھ لَیندے آں۔ ایس طرح اک-جگہ والے بنیادی
بازگشتی فنکشناں دی گنتی وار فہرست بناون دا اک
واضح طریقہ مل گیا اے۔

(اصل وچ، کجھ فنکشن، جداں مستقل صفر فنکشن، فہرست
وچ اک توں ودھ وار آون گے۔ ایہہ صرف ساڈی کوڈکاری
دی وجہ توں نہیں، بلکہ ایس وجہ توں وی اے کہ
مستقل صفر فنکشن دی اک توں ودھ علامتی لکھت اے۔
اگے اسی ویکھاں گے کہ ایہناں دہرائیاں توں قابلِ
حساب طریقے نال بچیا نہیں جا سکدا؛ مثال لئی، کوئی
قابلِ حساب فنکشن نہیں جیہڑا فیصلہ کرے کہ دِتّی
علامتی لکھت مستقل صفر فنکشن نوں ظاہر کردی اے
یا نہیں۔)

ہُن اسی فنکشن $g(x,y)$ نوں $f_x(y)$ نال
متعین کر سکدے آں، جتھے $f_x$ اوہی گنتی وار
فہرست دا حوالہ اے جیہڑی ہُن دسی اے۔ سانوں کِویں
پتا اے کہ $g(x,y)$ قابلِ حساب اے؟ بدیہی طور اُتے
ایہہ صاف اے: $g(x,y)$ حساب کرن لئی، پہلاں
$x$ دا کوڈ «کھولو» تے ویکھو کہ ایہہ اک-جگہ
والے فنکشن دی علامتی لکھت اے یا نہیں۔ جے اے،
تاں اوہ فنکشن دی قیمت انپٹ~$y$ اُتے حساب کرو۔

\begin{tagblock}{TMs}
\begin{digress}
ممکن اے تسی پہلاں ہی قائل ہوو کہ (کجھ محنت نال!)
ایہہ کم کرن والا پروگرام (مثال لئی، Java یا C++
وچ) لکھیا جا سکدا اے؛ تے ہُن اسی چرچ--ٹیورنگ
دعویٰ (Church--Turing thesis) دا سہارا لے سکدے
آں، جیہڑا کہندا اے کہ بدیہی سمجھ مطابق جیہڑی وی
شے قابلِ حساب اے، اوہ ٹیورنگ مشین نال حساب کیتی
جا سکدی اے۔

بے شک، $g(x,y)$ دے قابلِ حساب ہون نوں دکھاون
دا ہور سِدھا طریقہ ایہہ اے کہ اوہنوں حساب کرن
والی ٹیورنگ مشین صاف صاف بیان کیتی جاوے۔ خاص
طور اُتے ایس نال چرچ--ٹیورنگ دعوے تے بدیہی سمجھ
دا سہارا نہیں لینا پئے گا۔ چھیتی ہی ساڈے کول
اینیاں بناوتاں ہو جان گیاں کہ حساب دے اک نمونے
دا سہارا لے کے $g(x,y)$ دے قابلِ حساب ہون نوں
دکھا سکّیے؛ اوہ نمونہ ایسا اے جیہدے حسابی عمل دی ٹیورنگ مشین اُتے
\emph{نقل کیتی} جا سکدی اے: یعنی بازگشتی فنکشن۔
\end{digress}
\end{tagblock}

\end{document}
